\cotefile{104}
\scanpage{9}
\mspage{2}
i.e. on a, pour \(\deleted{\unsure{x}}\xi\in J_n(n')\) i.e. \(x\in J_n,\ \tau_n(x)=n'\)
\[
 u_\xi(d_{n'})=\xi .
\]
\deleted{De} Plus, \unsure{fort|faut|fut} encore, on a : \inserted{\(u_\xi\) induit un isom} \(J_{n'}\xrightarrow{\overset{u_\xi}{\sim}}(J_n)_{\leq\xi}\)

Lorsque \(n'=n\), on a \(J_{n,n}=\{d_n\}\), et
\[
 u_{d_n}=\mathrm{id}_{J_n}
\]

La connaissance de \(\mathbb N,\ (J_n)_{n\in\mathbb N},\ (\tau_n)_{n\in\mathbb N},\
(\rho_{n,n'})_{\substack{(n',\deleted{\unsure{n'|m}}n)\in\mathbb N\times\mathbb N\\ n'\leq n}}\) permet de reconstituer \deleted{à peu de \unsure{chose|choses} près} \deleted{essentiellement \unsure{la|le} sous-catégorie pleine}\inserted{la catégorie} \(M\), à \deleted{il} \unsure{équiv|équiva} près, il manque encore une \deleted{\unsure{description|descripteur}} \inserted{condition} pour composer les hom ``standart entre les \(J_n\), i.e ceux provenant des \(J_{n',n'}\). \deleted{\unsure{Donc|Don}} On considère \unsure{leurs|une} situation
\diagram{\begin{tikzcd}[ampersand replacement=\&,column sep=.8em,row sep=small]
 \& \& \eta \& \& \xi\\[-.8em]
 \& \& \rotatebox{-90}{\(\in\)} \& \& \rotatebox{-90}{\(\in\)}\\
 J_{n''}\arrow[rr,"u^{n'}_{\eta}"{pos=.3}] \& \& J_{n'}\arrow[rr,"u_\xi^n"{pos=.25,yshift=-.7ex}] \& \& J_n\\
 n'' \& \leq \& \unsure{n'|m'} \& \leq \& n
\end{tikzcd}}
on doit avoir
\[
 u_\xi^n\,u_\eta^{n'}=u_{\deleted{\unsure{\xi|a}}\zeta}^{n}
 \qquad\text{où}\quad \zeta=u_\xi^n(\eta)\overset{\unsure{\cdot|-}}{=}
 (u_\xi^n\,u_\eta^{n'})(d_{n''})
\]
qui exprime la \deleted{\unsure{composition|stabilité|compatibilité|stabilis}} \unsure{stb.lté} de la notion d'hom admissible \unsure{pour les|entre} \(J_n\) par composition,

\deleted{Il est évident que la \unsure{catég|caté}} \inserted{Soit \(M_0\) la catégorie définie par}
\[
 \mathrm{Ob}\,M_0=\mathbb N,\quad \mathrm{Hom}_{M_0}(n',n)=
 \left\{\begin{aligned}
 &\varnothing\ \text{si }n'\nleq n\\
 &\text{ens. des hom. admissibles}\\
 &\qquad J_{n'}\to J_n\quad\text{Si }n'\leq n
 \end{aligned}\right.
\]
% archive sheet number: 8

\scanpage{10}
la composition des $\operatorname{Hom}$ étant celle des $\mathrm{hom}$ \inserted{admiss.} entre les $ {\unsure{J|I}}_n$ (qui sont des immersions fermées). On a un foncteur canonique \deleted{\unsure{(défini|(définir}} \inserted{(indépend. \unsure{de de|de} tt choix, \deleted{\unsure{à|s}} si on y tient)}
\[
M\longrightarrow M_0
\]
qui est une équivalence de catégories -- en fait, \deleted{\unsure{cela|cel|et}} c'est un isom. de catégories, \unsure{ssi|NB} \deleted{\unsure{seulement}} $M$ est une catégorie réduite (deux objets isom sont identiques)

\deleted{\unsure{Je|La}} Je vais supposer pr la suite $M$ \unsure{réduite|réduit}. on peut donc l'identifier : la catégorie $M_0$, definie pr les données
\[
\resizebox{\linewidth}{!}{$\displaystyle
\left\{\begin{array}{@{}l@{\qquad}l@{\qquad}l@{}}
\mathbb N & (J_n)_{\in\mathbb N} &
,\;(\tau_n:J_n\to\mathbb N)_{n\in\mathbb N}\\
\text{ens ordonné} & \text{famille d'ens} &
\text{famille d'applict}\\
&\text{ordonnés, chaque }J_n\text{ \unsure{a} un}&\text{strictement}\\
&\text{plus grd élément}&\text{croissantes}\\
&&\tau_n(J_n)=\unsure{\mathbb N'|\mathbb N}_{\leq n}\\
\multicolumn{2}{@{}l}{(\rho_{n,n'})_{\substack{(n',n')\in\mathbb N\times\mathbb N\\n'\leq n}}}
&\tau_n^{-1}(n)=d_n\\
\multicolumn{2}{@{}l}{\text{famille d'applications}}&\text{(plus grd él. de ?}\\[1ex]
\multicolumn{3}{@{}l}{\begin{aligned}
\rho_{n,n'}:\tau_n^{-1}(n')&\longrightarrow\operatorname{Hom}(J_{n'},J_n)\\
\xi&\longmapsto (u^n_\xi:J_{n'}-J_n)
\end{aligned}}\\
\multicolumn{3}{@{}l}{\text{satisfaisant }\left\{\begin{aligned}
u^n_\xi &: J_{n'}\xrightarrow{\sim}(J_n)_{\leq\xi}&&
\text{(d'où }u^n_\xi(d_{n'})=\xi\text{)}\\
u^n_{\deleted{\unsure{\xi|\zeta|\eta}}}\circ u^{n'}_\eta&=u^n_\zeta&&
\text{\unsure{avec|où} }\zeta=u^n_\xi(\eta)
\end{aligned}\right.}\\[-.4ex]
\multicolumn{3}{r}{\text{pour }\xi\in J_n,\ \eta\in J_{n'},\
n'=\tau_n(\xi),\ n''=\unsure{\tau|E}_{n'}(\eta)}
\end{array}\right.
$}
\]
% archive sheet number: 9

\scanpage{11}
\mspage{3}
Inversement, de telles données définissent une catégorie $M_0$, satisfaisant $(i)(ii)$.

\uline{NB} Dans les cas que j'ai en tête, $\mathbb N=\mathbb N$

\uline{Ex 1} Catégorie des ens finis tot. ordonnés (ou : des \uline{simplexes} types $\Delta_n=\inserted{[0,n]}$) et des applictions strictement croissantes
% slanted
\margin{$J_n={\unsure{\mathfrak P|\mathscr P}}^{*}(\Delta_n)$\\
$\operatorname{Card} J_n-2^{n+1}-1$}

\uline{Ex 2} Catégorie des \deleted{cubes} \inserted{``\uline{cubes}} épinglés'', \deleted{i.e.} dont leurs objets \unsure{s\textsuperscript{t}|sont} les \deleted{\unsure{systèmes|applications} \unsure{fléchés|fléchées}} cubes combinatoires,
\diagram{\begin{tikzcd}[ampersand replacement=\&, row sep=small,
execute at end picture={\draw[thin] (h104arrbot.south west) -- (h104arrtop.north east);}]
|[alias=h104arrtop]|\deleted{\widetilde A}\arrow[d] \\
|[alias=h104arrbot]|\deleted{A}
\end{tikzcd}}
\deleted{$A,\widetilde A$ ens finis non vides} \deleted{$\widetilde A\unsure{-|=}A$} épinglés par un ordre total sur l'ens des paires d'hyperfaces, \deleted{ou} et le choix d'une hyper\unsure{f.|face}, ou face notée $+$ dans chaque paire.
% slanted
\margin{$J_n={}$ens des sections \unsure{p.tielles|partielles} de $\unsure{\{|(}\,\Delta_{n-1}\times\{\pm1\}\to\Delta_{n-1}\unsure{\}|)}$ avec \unsure{relat\textsuperscript{n}|relat\textsuperscript{o}} d'ordre \uline{opposée} de celle de prolongement\\
$\operatorname{Card} J_n=3^n\deleted{\unsure{4}\,\unsure{2^n}\,\unsure{-2^n}\,\unsure{(2^{n+1}-1)}}$}

T. épinglage du cube induit un épinglage de toute facette, donc ... \deleted{\unsure{I|J|l}}

\uline{Ex. 3} Catégorie \deleted{des} \inserted{des} \uline{hémisphèr\unsure{es}}\deleted{\unsure{iques}} $D_n$, où $J_n=(\Delta_n\times\{\pm1\})\amalg\{d_n\}$, avec $\{\pm1\}$ muni de l'ordre chaotique, $\deleted{\{}d_n\deleted{\}}$ le plus grand élément de $J_n$.
% slanted
\margin{$\operatorname{Card} J_n=2n+1$}

Dans \deleted{tous} ces \inserted{trois} exemples, on a

\textcircled{a} \inserted{Si $n\geq1$} La réalisation top. $|\deleted{\unsure{le}} J_n^*|$ de $J_n^*\inserted{=J_n\setminus\{d_n\}}$ est une \deleted{\unsure{catég|catégorie}}

\deleted{\textcircled{b} \unsure{Sa|La|les} \unsure{réalisation|réalisations}} $(n-1)$ sphère, (donc $|J_n|$ est un $n$-disque)
% archive sheet number: 10

\scanpage{12}
\textcircled{b} Les ens. \inserted{ordonnés} $J_n^*$, (regardés comme un $(n-1)$-polyèdre, sont \uline{réguliers}, i.e. le groupe des autom est transitif sur l'ens des repères de $J_n^*$, i.e. l'ens des \unsure{chaines|chaînes} maximales
% slanted
\margin{b) est-il conséquence de a) ??}
\[
{\deleted{\unsure{x}}\unsure{\xi|x}}_0\unsure{<|\leftarrow}{\deleted{\unsure{x}}\xi}_1<.\,-<{\deleted{\unsure{x}}\xi}_{n-1}
\]

\uline{Question} Y a-t-il d'autres exemples d'une suite \deleted{$(P_n\unsure{\in|(\!-|=}J_{n+1}^{*})$} $ \deleted{(}\underset{\substack{\llcorner\mkern-5mu\underline{\phantom{J_n^*}}\mkern-5mu\lrcorner\\\inserted{P_{n-1}}}}{\deleted{J_n^*}}\deleted{)_{n\in\mathbb N,n\geq1}}$ \inserted{$(P_n)_{n\geq1}$} d'ens ordonnés \unsure{présentant|possédant} les propriétés
% slanted
\margin{$P_n=J_{n+1}^{*}$}

a) $P_n$ est un \inserted{$n$-}polyèdre \uline{\unsure{régulier|réguliers}} sphérique

b) Si $\xi\in P_n$, \deleted{\unsure{est de dim., si|est de dim. $n'$}} \inserted{alors} $(P_n)_{\leq\xi}$ \deleted{est un polyèdre combinatoire de dimension $n'$, \unsure{on} a} $(P_n)_{\leq\xi}\simeq P_{n'}$, pour $n'$ convenable $\leq n$.

(\unsure{Ça|Sa} doit etre dans les livres de Coxeter !)

On aimerait de plus, surtout, être capables de trouver des $\deleted{\unsure{\rho_{\xi,n}|u^n_\xi}}\,\rho_{n,n'}$ satisfaisant la condition de transitivité
% archive sheet number: 11

\scanpage{13}
\mspage{4}
Je considère maintenant la sous catégorie ${\unsure{\mathscr B|\mathcal B|B}}_0$ de $\widehat M\supset M$, formée des objets $X$ de la forme suivante :
\[
\left(\begin{aligned}
&X\text{ est }\deleted{\unsure{\varinjlim|\underset{\longrightarrow}{\text{limite}}}}\inserted{\text{sup}}\text{ de ses sous-objets isomorphes à des}\\
&\deleted{\text{\unsure{les|ses|des}}}\qquad\text{objets de }M\text{ (modèles)}
\end{aligned}\right.
\]
% slanted
\margin{Cela \unsure{implique|signifie} que \unsure{tt|$\forall$} hom\unsure{,} \deleted{de} $S\unsure{-|\to}X$, où $S\in\operatorname{Ob}M$, est un mono}

Si $K,L\subset X$ sont des sous-objets modèles, alors $K\cap L$ est sup \unsure{des|de} sous-objets de $K\cap L$ qui sont des modèles

Soit Donc $\widetilde X$ l'ens des \deleted{\unsure{obj}} sous-objets de $X$ qui sont des modèles. Alors $\widetilde X$. est un ens ordonné \inserted{$K$}, muni de la structure suivante \inserted{(``$M$-épinglage'' de $K$)} : \deleted{\unsure{à|l'} \unsure{une|savoir}}

a) \deleted{\unsure{un|une} \unsure{de} \unsure{Pour}} Une application strictement croissante
\[
K\xrightarrow{\tau_K}\deleted{\unsure{I|N}}N
\qquad\text{(application ``\uline{type}'')}
\]

b) Pour tt $\xi\in K$, \inserted{posant $n=\tau_K(x)$,} une \deleted{\unsure{applic\textsuperscript{n}|applic\textsuperscript{t}|applic}} \unsure{isomorphisme|isomorphis\textsuperscript{me}}
\[
\deleted{\unsure{D|I|\varphi_\xi}}\;J_n\underset{u_\xi^K}{\simeq}K_{\leq\xi}
\qquad\text{(application d'épinglage de }K_{\leq\unsure{\mid|\xi|\eta}}\text{)}
\]
% slanted
\margin{les applications $J_n\unsure{-|\to}K$ définies par les $\xi\in K(n)=\unsure{(}\unsure{\mathrm{d.}}\tau^{-1}(n)$ sont dite \uline{admissibles}}
\[
u_\xi^K(d_n)=\xi\quad\text{(= plus grd \unsure{elément|élément|elém|élém} de }K_{\leq\xi}\text{)}
\]
avec \deleted{condition} les propriétés \deleted{de \unsure{transitivité|transit|transit.}} :
% archive sheet number: 12

\scanpage{14}
\begingroup\makeatletter\tagsleft@true\makeatother
\[
\tau_K\circ u_\xi^K=\tau_n\tag*{(1)}
\]
et la propriété de transitivité pour les applications admissibles
\[
\begin{aligned}[t]
&\forall\,\xi\in K,\ \eta\in J_n\quad(n=\tau_K(\xi))\ \text{on a, si }n'=\tau_n(\eta)\\
&\qquad\deleted{\unsure{\tau|u}}\,u_\xi^K\circ u_\eta^n=u_\zeta^K,\quad\text{où }\zeta=u_\xi^K(\eta),
\end{aligned}\tag*{(2)}
\]
\endgroup
enfin

\deleted{(3) Si $\xi\in K,\ \eta\in K$, alors}

La catégorie dont les objets sont les objets admissibles de $\widehat M$, les flèches les \unsure{Isom|isom|mono|hom} entre eux, est équivalente : celle des ens ordonnés $M$-épinglés.

D'autre part, les hom quelconques \inserted{dans $\widehat C$ entre tels objets} se réconstituent également, ce sont les applictions \uline{strictement} \uline{croissantes}
\[
K\xrightarrow{f}L
\]
telles que $\left\{\begin{aligned}
&\text{a) }\tau_L\circ f=\tau_K\\
&\text{b) pour tt }\xi\in K,\ \inserted{\text{posant }\eta=f(\xi)\in L,}\ \text{on ait}
\end{aligned}\right.$ commutativité \unsure{du|de}
\diagram{\begin{tikzcd}[ampersand replacement=\&,column sep=small,row sep=large]
K_{\leq\xi}\arrow[rr,"{f|K_{\leq\xi}}"] \& \& {\deleted{K}L}_{\leq\eta}\\
\& J_n\arrow[ul,"\sim"']\arrow[ur,"\sim"] \& \& \text{où }n=\tau_K(\xi)=\tau_{\deleted{\unsure{K|\eta}}L}(\eta),
\end{tikzcd}}
% archive sheet number: 13

\scanpage{15}\mspage{5}
ce qui implique en particulier que $f\mid K_{\leq\unsure{\xi}}$ induit un iso
\[
K_{\leq\xi}\xrightarrow{\sim}{\deleted{\unsure{K|V}}L}_{\leq\eta}
\]

Nous nous intéressons surtout au cas où $f$ est injectif, ce qui correspond au \deleted{fait \unsure{que}} cas \unsure{d'un|d'une} \uline{mono} dans $\widehat C$.

Alors $\unsure{f|f'}$ induit un \uline{isom}. de $K$ sur \unsure{\deleted{un}|un} \deleted{sous-ens} son image $L'=f(K)$, muni de l'ordre induit par $L$. Notons que $L'$ est une partie \uline{fermée} de $L$ $(x\in L',y\leq x\Rightarrow y\unsure{\in}L,$ et que tte partie fermée admet un $M$-épinglage \deleted{induit} unique, tel que l'inclusion soit admissible (épinglage induit) -- et revenant : $f$, bien sûr $K\overset{f}{\xrightarrow{\sim}}L'$, est compatible avec les $M$-épinglages.
% archive sheet number: 14

\scanpage{16}
Un cas intéressant est celui où $\deleted{\unsure{K|X}}\unsure{X|K}$ \deleted{\unsure{admet}} \inserted{est \unsure{le|la} sup de un nb fini de sous-}\deleted{un nb fini} de \deleted{\unsure{pts} maximaux, et} objets modèles, ce qui signifie que l'ens des éléments maximaux de $K\inserted{=\unsure{\widetilde X|X}}$ est fini, et que tt élément de $K$ est majoré par un élément maximal. Pour un ens. ordonné ayant \deleted{\unsure{ces}} \inserted{la $2^{\text{\unsure{ième|iem|ème}}}$} \unsure{propriété|propriétés}, un $M$-épinglage est connu quand on connait les $u^K_{\deleted{\xi}\alpha}$ pour $\deleted{\xi}\alpha\inserted{\in K}$\unsure{(|,} maximal.

La condition pour que des donn.
\[
u^K_{\deleted{\xi}\alpha}: J_{\tau_n(\deleted{\xi}\alpha)}
   \xrightarrow{\sim} K_{\leq\deleted{\xi}\alpha}\qquad
   (\unsure{\xi}\in\operatorname{Max}K,
\]
définissent un épinglage, est que pour \deleted{\unsure{tt|$\xi$|$\forall$}} $\alpha,\beta\in\operatorname{Max}K$ et $\xi\in K$, $\xi\leq\alpha,\beta$, \deleted{posant} $n=\tau_K(\xi)$, \unsure{les|des} deux iso
\[
J_n\xrightarrow{\sim}K_{\leq\xi}
\]
définis par l'épinglage \deleted{\unsure{de|des}} ${\deleted{\unsure{\tau}}u}^K_\alpha$ de $K_{\leq\alpha}$, et par l'épinglage $u^{\unsure{K|k|n}}_\beta$ de $K_{\leq\beta}$, \deleted{\unsure{soient}} les \unsure{mêmes}.
% archive sheet number: 15

\scanpage{17}\mspage{6}
La catégorie des objets admissibles de $\widehat C$ est stable par \deleted{\unsure{les} l'opération} sommes quelconques, et par sommes amalgamées $X\amalg_Z Y$, \unsure{don|dont} $Z\unsure{\mathord{-}|\to}X$ et $Z\mathord{-}Y$ des mono, et le foncteur $X\to\widetilde X$ commute : ces opérations, dans un sens évident. De plus, les objets de type fini donnent des objets de t.f., pourvu, \unsure{(}dans le cas des sommes\unsure{)|,}, qu'on se borne \unsure{:|à} des sommes \uline{finies}.

\uline{NB} Quand les $\unsure{J_n|I_n|J}$ sont finis, alors les objets \inserted{adm.} de t.f. de $\widehat C$ \deleted{\unsure{aussi|sont} \unsure{(par ex)|t.f. les $X$ t.|de les $X$,}} \inserted{correspondent aux} ens ordonnés $M$-épinglés $K$ qui \deleted{sont} \uline{finis}.
% archive sheet number: 16
